% !TEX root = ../main.tex
%%%%%%%%%%%%%%%%%%%%%%%%%%%%%%%%%%%%%%%%%%%%%%%%%%%%%%%%%%%%%%%%%%%%%%%
% Results III - fiscal conditionality
%%%%%%%%%%%%%%%%%%%%%%%%%%%%%%%%%%%%%%%%%%%%%%%%%%%%%%%%%%%%%%%%%%%%%%%

\section{Results III: the sign region is conditional on fiscal institutions}
\label{sec:government}

The results so far describe an economy with no public sector, in which the
unemployed earn nothing and therefore leak entirely out of the circular flow.
That is a modelling choice, not a neutral baseline, and this section shows it
carries the sign structure of \cref{sec:shape} with it. We reinstate a minimal
government: a flat tax on accrued income funds an equal transfer to the
unemployed, indexed to the current wage, with the budget balanced every
period. One economic parameter is added, the replacement rate $rr$, whose
range is anchored to OECD net replacement rates (\cref{tab:params}).

\paragraph{Dose-response: crowding-in.} At the headline demand-constrained
scenario the benefit lowers unemployment and raises \emph{both} output and
capital (\cref{tab:government}).

\begin{table}[htbp]
\centering
\small
\caption{Fiscal dose-response at $c_0 = 1.0$, $\sigma = 0.5$, $\eta = 0.10$,
$\rho = 0.40$.}
\label{tab:government}
\begin{tabular}{lrrrrrr}
\toprule
$rr$ & $U$ & $Y$ & $K$ & Wage share & Realised tax & Cash-constrained\\
\midrule
0.00 & 0.566 & 82.1 & 298.8 & 0.440 & 0.000 & 0.90\\
0.25 & 0.483 & 98.1 & 359.2 & 0.446 & 0.128 & 0.90\\
0.50 & 0.427 & 108.7 & 396.6 & 0.452 & 0.206 & 0.90\\
0.75 & 0.373 & 119.0 & 435.7 & 0.457 & 0.250 & 0.90\\
\bottomrule
\end{tabular}
\end{table}

Capital rises from 299 to 436. In a demand-constrained regime redistribution
is \emph{crowding-in}, because more demand means more profit and hence more
investment through $I = \rho\Pi$ --- the opposite of the crowding-out a
loanable-funds account would predict, and a direct consequence of financing
investment out of retained earnings rather than out of saving. The
cash-constrained fraction stays at $0.90$, all ninety workers, invariant to
$rr$: the benefit never lifts a worker off the liquidity constraint, which is
precisely why the balanced-budget multiplier keeps delivering across the whole
dose.

\paragraph{The benefit eliminates the wage-led region.} At $rr = 0.5$, inside
the empirically anchored band, $\sigstar$ becomes \emph{undefined}: every
estimated $\mathrm{d}Y/\mathrm{d}\rho$ turns positive across the tested range
($+38.7$ at $\sigma = 1$; $+19.3$ at $\sigma = 1.5$), pushing $\sigstar$ above
$1.5$. The demand floor makes retention expansionary at \emph{every}
elasticity tested. The unemployment frontier, by contrast, barely moves.

\begin{declbox}
The region in which retention depresses output is not a property of the
technology alone. It exists because the unemployed leak out of the circular
flow, and a balanced-budget benefit at a replacement rate inside the OECD band
closes that leak and removes the region. \textbf{The sign is conditional on
fiscal institutions}, and any statement of it that omits the institutional
setting is incomplete.
\end{declbox}

\paragraph{The same instrument fails on the collapse.} The falsifiable guess
was that a demand floor would also shrink the collapse region of
\cref{sec:wagecurve}. It does the opposite: cells with any collapsed seed rise
from 16 to 26 ($\eta = 0.10$) and 16 to 29 ($\eta = 0.15$), and the mean
fraction of seeds at $U = 1$ roughly doubles. In the reference cell the economy
dies at both $rr = 0$ and $rr = 0.5$, but at $rr = 0.5$ the tax is \emph{pinned
at its cap} in 100\% of periods --- the instrument saturates and the economy
still dies. Three mechanisms explain it: where firms are collapsing the tax
base is almost entirely wages, so taxing workers to pay workers is
MPC-neutral; the benefit is indexed to $w_t$ and is therefore
\emph{procyclical}, weakest exactly when it is most needed; and the demand it
does inject \emph{amplifies} the wage$\,\to U \to\,$capital-erosion
oscillation. That the same instrument can close the leak that produces the
wage-led region and yet enlarge the collapse is the sharpest available
statement of the boundary drawn in \cref{sec:discussion}: it acts on demand,
and only one of the two phenomena is a demand phenomenon. \Cref{fig:government}
shows the two faces side by side.

\begin{figure}[htbp]
\centering
\begin{subfigure}{0.49\textwidth}
  \includegraphics[width=\linewidth]{ces_b09_dose_response.png}
  \caption{Dose-response: $U$, $Y$, $K$ against $rr$.}
\end{subfigure}\hfill
\begin{subfigure}{0.49\textwidth}
  \includegraphics[width=\linewidth]{ces_b09_collapse_map.png}
  \caption{Collapse map: the demand floor enlarges it.}
\end{subfigure}
\caption{The two faces of the balanced-budget benefit.}
\label{fig:government}
\end{figure}
